\subsection{p 进李群}

这样的群最早于 1907 年出现在亨塞尔关于 p 进解析函数（由整级数展开定义）的工作{[}157 e{]}中。他特别研究了指数函数与对数函数；尽管定义它们的级数的行为先验地看来令人惊讶（例如指数级数并非处处收敛），它们的函数性质仍然成立，这就给出了 \(Q_{p}\) 的加法群与乘法群之间的一个局部同构（或者更一般地，给出特征为零的每个完备超度量域的加法群与乘法群之间的局部同构）。

A. 韦伊{[}330 a{]}与 E. 卢茨{[}210{]}关于 p 进椭圆曲线（1936 年）的工作，所关心的同样是交换群（不过这一次不是线性的）。除算术应用之外，其中还可以找到该群与加法群之间一个局部同构的构造，它建立在一个不变微分形式的积分之上。这一方法同样可以应用于阿贝尔簇，C. 沙博蒂不久之后指出了这一点，他未作更多说明便使用这一方法证明了“莫德尔猜想”的一个特例{[}61{]}。

从这时起人们已经清楚，李群的局部理论几乎可以不加改变地应用于 p 进情形。李群—李代数“词典”的基本定理是 1942 年在 R. 胡克（谢瓦莱的学生）的博士论文{[}166{]}中建立的；这项工作还包含关于实李群闭子群的 E. 嘉当定理的 p 进类似物。

更近些时候，M. 拉扎尔{[}195 b{]}对 \(\mathbf{Q}_p\) 上的紧解析群发展了“词典”的一种更精确的形式。他证明，\(p\) 进解析结构在紧群 \(G\) 上的存在与 \(G\) 上某些过滤的存在紧密相连，并给出了它的多种应用（例如对 \(G\) 的上同调）。拉扎尔的工具之一是对邓金关于 \(p\) 进豪斯多夫级数收敛性结果的一个改进{[}195 a{]}。

\section{自由李代数}

我们还须谈谈一系列关于李代数的工作，它们与李群理论的联系十分微弱；但另一方面，这些研究在“抽象”群理论中、尤其是在幂零群理论中有着重要的应用。

这件事的起源是 P. 霍尔发表于 1932 年的工作{[}144{]}。不过那里完全谈不上李代数：P. 霍尔着眼于研究某一类 p 群，即他称为“正则”的那类 p 群。但这导致他仔细考察了一个群的迭代换位子和下中心列；他借此机会建立了雅可比恒等式的一个变体，以及“霍尔公式”

\[
(x y) ^ {n} = x ^ {n} y ^ {n} (x, y) ^ {n (n - 1) / 2} \dots
\]

不久之后（1935—1937 年），W. 马格努斯{[}215 a 和 b{]}与 E. 维特{[}337 b{]}的基本工作问世。在{[}215 a{]}中，马格努斯使用了与豪斯多夫相同的形式级数代数 \(\hat{A}\)（后来称为“马格努斯代数”）；他把自由群 \(F\) 嵌入其中，并利用 \(\hat{A}\) 的自然过滤，得到了一个递降序列 \((F_n)\)，其中的群都是 \(F\) 的子群；这是最早的过滤例子之一。他猜想 \(F_n\) 与 \(F\) 的下中心列的各项相重合。这一猜想在他的第二篇论文{[}215 b{]}中得到证明；也正是在那里，他明确地把他的思想与 P. 霍尔的思想联系起来，并且定义了自由李代数 \(L\)（作为 \(A\) 的一个子代数），而且实质上证明了它与分次的 \(F\) 相等同。在{[}337 b{]}中，维特在若干方面完善了这一结果。他特别证明了 \(L\) 的包络代数是一个自由结合代数，并由此立即导出 \(L\) 的齐次分量的秩（“维特公式”）。

至于以“霍尔基”之名著称的 L 的基的确定，看来它直到 1950 年才第一次出现，即出现在 M. 霍尔的一篇短文{[}143{]}中，尽管它已隐含在上文引述的 P. 霍尔与 W. 马格努斯的工作里。

\chapter{由反射生成的群；根系。}

本注所考虑的群，是同几何学、分析学以及李群理论的各种问题相关联地出现的，它们有时以置换群的形式出现，有时以欧几里得几何或双曲几何中位移群的形式出现，而这些不同的观点直到晚近才得到协调。

从历史上看，这一理论的开端要比群概念的引入早得多：它的源头其实在于对几何图形的“正则性”或“对称性”的研究，尤其是正多边形与正多面体的确定（无疑可以追溯到毕达哥拉斯学派），后者构成了欧几里得《几何原本》的压卷之作，也是希腊天才最令人赞叹的创造之一。后来，特别是在中世纪盛期的阿拉伯作者那里，然后在开普勒那里，出现了用两两全等（但不必正则）的多边形对平面或球面作正则“铺砌”的数学理论的雏形，这无疑与古代文明和阿拉伯文明所构想的各种类型的装饰的起源有联系（这些装饰有充分的理由被视为这些文明所发展的数学的一个真实组成部分{[}291{]}）。

大约在 1830 至 1840 年间，晶体学方面的研究（赫塞尔、布拉维、默比乌斯）引导人们考虑一个问题，它恰好就是确定三维欧几里得空间中有限位移群的问题，尽管上述作者当时还没有使用群论的语言；群论的语言直到 1860 年前后才逐渐通行，而正是在对群进行分类的背景之下，若尔当于 1869 年{[}174 b{]}确定了 \(\mathbf{R}^3\) 中保持定向的离散位移子群（更一般地，确定了保持定向的位移群的所有闭子群）。

直到十九世纪的最后几年，这股思想潮流都在朝许多方向发展，其中最显著的如下：

\begin{enumerate}
\def\labelenumi{\arabic{enumi}.}
\tightlist
\item
  与有限群理论中出现得很早的一种倾向相一致，人们寻求用简单类型的生成元和关系来给出有限位移群的“表出”。正是这样，哈密顿早在 1856 年{[}145 b{]}就证明了欧几里得空间中
\end{enumerate}

\(R^{3}\) 的有限旋转群由两个生成元 S、T 生成，它们满足关系 \(S^{p} = T^{q} = (ST)^{3} = 1\)，其中 p 与 q 取适当的值。

\begin{enumerate}
\def\labelenumi{\arabic{enumi}.}
\setcounter{enumi}{1}
\item
  离散位移群可以含有反射，也可以不含有。早在 1852 年，默比乌斯就实质性地确定了球面几何中由反射生成的有限位移群（这等价于 \(\mathbf{R}^3\) 中有限欧几里得位移群的同一问题）；他发现，除去循环群之外，这样的群都以一个球面三角形为基本域，其角具有形式 \(\pi / p, \pi / q, \pi / r\)，其中 \(p, q, r\) 是三个 \(>1\) 的整数，满足 \(\frac{1}{p} + \frac{1}{q} + \frac{1}{r} > 1\)（{[}223{]}, v. II, pp.~349-360 与 pp.~561-708）。他还注意到，这些群作为子群包含着所有的有限位移群。
\item
  当继黎曼和施瓦茨关于超几何函数与共形表示的工作之后，开始研究用以为圆弧所界的图形对复平面或半平面所作的“铺砌”时，上述最后这股思想潮流获得了新的扩展；克莱因和庞加莱把它奠立为“自守函数”理论的基础，并在其中（就圆弧与一条固定直线正交的情形）认识到一个与求非欧双曲平面（等同于“庞加莱半平面”）位移群的离散子群等价的问题{[}118{]}。
\item
  正多面体的概念以及用这种多面体对 \(R^{3}\) 作铺砌的概念，被施莱夫利推广到了所有的欧几里得空间 \(R^{n}\)，这项工作可追溯到 1850 年前后，但很久以后才得以发表，并且在很长时间里一直无人问津（{[}273{]}, v. I, pp.~167-387）；他完全确定了每个 \(R^{n}\) 中的正“多胞形”、保持这种多胞形不变的位移群，以及该群的一个基本域，与默比乌斯研究过的 n = 3 的情形一样，这个基本域是一个“房间”，它在球面 \(S_{n-1}\) 上的迹是一个球面单形。然而，他没有着手处理求 \(R^{n}\) 中由反射生成的有限位移群这一逆问题；这个问题很久以后才得到解决：n = 4 的情形由古尔萨{[}132{]}解决，而对任意的 n，其解决则要等到 E. 嘉当（{[}52 a{]}, v. I₂, pp.~1003-1020）与考克斯特{[}70 c{]}的工作，我们稍后将再回到这些工作。
\end{enumerate}

1890 年前后，随着基灵和 E. 嘉当关于李群最初的工作，一股新的思想潮流开始了，它在很长时间里将与前面的那些工作互无联系地发展。基灵{[}180{]}和嘉当（{[}52 a{]}, v. I \(_{1}\), pp.~137-287）在研究复半单李代数的结构时，从一开始就打算让某些线性形式 \(\omega_{\alpha}\) 起着首要的作用，这些线性形式定义在这样一个李代数的一个“嘉当子代数”\(\mathfrak{h}\) 上，而这个李代数就是 \(\mathfrak{g}\)；它们就是相对于 \(\mathfrak{h}\) 的“根”，之所以如此称呼，是因为在基灵那里它们表现为特征方程 \(\det(\mathrm{ad}_{\mathfrak{g}}(x) - T) = 0\) 的根，这里把它们看作 \(x \in \mathfrak{h}\) 的函数。基灵和嘉当所建立的这些“根”的性质，用今天的几何语言来说，无非是断言它们构成一个“约化根系”；

他们随后证明，复半单李代数的分类可以归结为相伴“根系”的分类，而后者本身又可以归结为某些整系数矩阵（后来称为“嘉当矩阵”）的确定。基灵和嘉当还揭示出，对每个根 \(\omega_{\alpha}\)，存在根集合的一个对合置换 \(S_{\alpha}\)；他们本质性地使用了变换 \(C = S_{\alpha_{1}}S_{\alpha_{2}} \ldots S_{\alpha_{l}}\)，即与构成一个基础根系的 l 个根相伴的诸置换的乘积（这一变换今天称为“考克斯特变换”）；他们甚至把这个置换扩张为由基础根 \(\omega_{\alpha_{i}}(1 \leq i \leq l)\) 生成的向量空间上的一个线性变换，并研究其特征值（{[}180{]}, II, p.~20；{[}52 a{]}, v. I \(_{1}\), p.~58）。但无论是基灵，还是起初的嘉当，似乎都没有想到考虑群 \(G'\)，即由诸 \(S_{\alpha}\) 生成的群；而当嘉当稍后（{[}52 a{]}, v. I \(_{1}\), pp.~293-353）着手确定特征方程

\[
\det (\mathrm{ad} _ {\mathfrak {g}} (x) - T) = 0
\]

（即“一般元素”\(x \in h\) 的特征方程）的伽罗瓦群 G 时，他起初并没有援引 \(S_{\alpha}\) 来研究它；三十年后，已经处于 H. 外尔工作的影响之下，他证明了（{[}52 a{]}, v. I \(_{1}\), pp.~555-568）G 以群 \(G'\) 为正规子群，并在所有情形确定了商群 \(G/G'\) 的结构：对单代数 g 而言，它是 1 阶或 2 阶的，但 \(D_{4}\) 型除外，那时它同构于 \(S_{3}\)；也是在这次机会中，他把 \(G'\) 解释为复半单李代数中使一个嘉当子代数固定不动的内自同构所诱导的群。

我们刚才提到的 H. 外尔的工作，开创了群 \(G'\)（后来称为 g 的“外尔群”）的几何解释；正如基灵和嘉当对变换 C 所做的那样，他想到把 \(S_{\alpha}\) 看作 h 上线性形式向量空间中的反射。也是在 H. 外尔的论文{[}331 b{]}中，可以看到“仿射外尔群”的基本域的出现（不过它与“外尔群”\(G'\) 的联系并没有被明确指出）；外尔利用它证明了紧半单群的基本群是有限的，这是他证明复半单李代数线性表示完全可约性的一个关键点。不久之后，E. 嘉当完成了综合，把 H. 外尔的整体观点、他自己的实或复半单李代数理论以及他当时正在建立的对称黎曼空间理论融为一体。在论文（{[}52 a{]}, v. I₂, pp.~793-840）中，他完成了外尔群与仿射外尔群的基本多胞形的确定，并引入了权系和根基权系；他把这一讨论推广到对称空间，从而特别遇到了非约化根系的第一批例子（{[}52 a{]}, v. I₂, pp.~1003-1010）。最后，论文（{[}52 a{]}, v. I₂, pp.~867-989）给出了下列事实的第一个证明：\(R^{n}\) 中每个由反射生成的不可约有限群都具有一个基本域，它在球面 \(S_{n-1}\) 上的迹是一个球面单形；也正是在这项工作中，他借助几何的考虑证明了最大根的唯一性（在根系上任取的一个字典序下）。

稍后，范德瓦尔登{[}317 b{]}依靠 H. 外尔的论文证明了复半单李代数的分类等价于约化根系的分类，并通过初等的几何考虑完成了这一分类（而在基灵和嘉当那里，这一分类是复杂的行列式计算的结果）。大约同一时期，考克斯特明确地确定了所有由反射生成的不可约有限欧几里得位移群{[}70 c{]}；他就这样补全了 E. 嘉当论文（{[}52 a{]}, v. I₂, pp.~793-840）的结果，因为嘉当当时只确定了“晶体”群（即与一个根系相伴的群，或者也可说能够嵌入一个无限离散位移群的群）。第二年{[}70 d{]}，考克斯特证明了：由反射生成的有限群是仅有的（在相差一个同构的意义下）具有如下表出的有限群——由对合生成元 \(R_i\) 加上形如 \((R_i R_j)^{m_{ij}} = 1\) 的关系（\(m_{ij}\) 为整数）给出；从那时起，凡是具有这种表出的群（有限的或否）都得到了“考克斯特群”这一名称。

我们在上面描述过的两股研究潮流之间的第一个联系，看来是由考克斯特{[}70 e{]}、随后由维特{[}337 c{]}建立的。他们注意到，由反射生成的不可约无限欧几里得位移群与复单李代数（在相差一个同构的意义下）一一对应。维特给出了这类离散群的一个新的确定，并且还推广了上面回顾的考克斯特定理{[}70 d{]}，同时刻画了同构于无限离散欧几里得位移群的那些考克斯特群。这一结果，加上双曲几何中的类似群也是考克斯特群这一事实，\(^{1}\) 使人们得以公开着手研究这些群，先是把重点放在几何实现上（{[}71{]}、{[}308 a{]}），随后按照 J. 蒂茨{[}308 b 和 c{]}的做法在纯代数的框架内进行。

从维特的工作开始，半单李群的理论与由反射生成的离散群的理论将不断地彼此发生作用，而且成果极其丰硕。早在 1941 年，施蒂费尔{[}296{]}就指出，外尔群恰好是使一个格保持不变的有限反射群。谢瓦莱{[}62 e{]}和哈里希-钱德拉{[}149 a{]}在 1948 至 1951 年间对“晶体”群与复半单李代数之间的双射对应给出了先验的证明；在那之前，人们除了对每一类型的单李代数分别验证这一对应之外，别无他法。

另一方面，大约在 1950 年，人们注意到在外尔群下不变的多项式在无穷维线性表示的理论{[}149 a{]}以及李群的拓扑学中起着重要的作用。

就考克斯特而言，他{[}70 g{]}重新拾起了对变换 C（由反射生成的有限群 W 的基础反射的乘积）的研究，并注意到（通过对每个类型分别考察）W 下不变的多项式代数由代数无关的元素生成，这些元素的次数与 C 的特征值以简单的方式相联系。这些结果的先验证明随后由谢瓦莱{[}62 f{]}就第一个结果、由科尔曼{[}68{]}和斯坦伯格{[}292{]}就第二个结果给出。

随着 A. 博雷尔关于线性代数群{[}30{]}的工作，李群理论中开始了一些新的发展，这些发展后来导致了这一理论的显著扩展。A. 博雷尔揭示出李群中极大连通可解子群（后来称为“博雷尔子群”）的重要性，并把它们当作主要工具，用来把经典理论的一大部分移植到代数闭域上的代数群（不过当时还没有得到单代数群的分类\(^{2}\)）。（在实或复经典群的情形）博雷尔子群几年之前就已经出现在盖尔范德和纽马克关于无穷维表示的工作中；1954 年，F. 布吕阿发现了一个引人注目的事实：对经典单群而言，群按博雷尔群所作的双陪集分解以典范的方式由外尔群来标号{[}40{]}。这一结果随后由哈里希-钱德拉{[}149 b{]}推广到所有的实和复半单群。另一方面，1955 年，谢瓦莱{[}62 g{]}成功地做到了：对每个复半单李代数 g 和每个交换域 k，联系一个系数在 k 中的矩阵群，它具有布吕阿分解；他利用后一事实证明，除去少数例外，这样定义的群是单群（在抽象群理论的意义下）。他就这样“解释”了最迟自若尔当和李以来就已注意到的重合，即 A、B、C、D 型的单李群（在李群理论的意义下）与在任意域上以纯代数方式定义的经典单群之间的重合（直到那时，这一重合只能被推广到例外型 \(G_{2}\) 和 \(E_{6}\)，那是迪克森做到的{[}88c 和 d{]}）。特别地，取 k 为有限域时，谢瓦莱的构造对每一类型的复单李代数都给出一族有限单群，其中包含当时已知的有限单群的很大一部分，还包含三个新的系列（对应于单李代数 \(F_{4}, E_{7}\) 和 \(E_{8}\) 的类型）。不久之后，多位作者（赫齐格、铃木、Ree、斯坦伯格和蒂茨）用各种程序——即对谢瓦莱方法的种种修改——一方面证明了可以用类似的方式得到当时已知的其他有限单群（交错群与马蒂厄群除外），另一方面构造出了新的有限单群的其他系列（参见{[}54{]}）。

几乎在同一时期，谢瓦莱{[}62 h{]}仍然使用布吕阿分解的技术，再加上关于博雷尔子群正规化子的一个关键结果，重新研究了线性代数群，并得到了这样的结果：在任意特征的特征代数闭域 k 上，半单线性代数群\(^{3}\)的理论本质上给出与 k = C 时的基灵—嘉当分类相同的类型。接着，J. 蒂茨{[}308 a 和 b{]}通过分析谢瓦莱的方法，得到了布吕阿分解的一个公理化版本（“BN 对”），其形式异常灵活，只涉及群的结构；目前正是这一概念以“蒂茨系”的名称为人所知。我们在上面谈到的所有单群（该词各种意义下的单群）都典范地带有蒂茨系，而蒂茨本人{[}308 c{]}证明了：在抽象群 G 中这样一个系的存在，加上来自纯群论的几个补充性质，就使人们能够证明 G 是单群，这条定理涵盖了迄今为止对这些群所给出的绝大多数单性证明。另一方面，他与 A. 博雷尔合作，推广了{[}62 h{]}中谢瓦莱的结果，证明了在任意域上的半单线性代数群的有理点群中存在蒂茨系{[}31{]}。

这些问题中出现的所有蒂茨系都具有有限的外尔群。另一类例子由岩堀和松本{[}170{]}发现；他们证明了，如果在{[}62 g{]}中谢瓦莱的构造里取 k 为一个 p 进域，那么所得到的群带有一个蒂茨系，其外尔群是出发时所用的那个复半单李代数的仿射外尔群。这一结果刚刚由布吕阿和蒂茨{[}41{]}推广到局部域上的所有半单代数群。

\specialchapter{参考文献}

1 N. H. Abel. Oeuvres, 2 vol.~ed.~Sylow 与 Lie, Christiana, 1881.

2a I. Ado. 论借助线性代换表示连续有限群（俄文）, Bull. Phys. Math. Soc. Kazan, v. VII (1935), pp.~3-43.

2b I. Ado. 李代数的矩阵表示（俄文）, Uspehi Mat. Nauk., v. II (1947), pp.~159-173（英文译本：Amer. Math. Soc. Transl., (1), vol.~9, pp.~308-327）.

3 A. D. Alexandroff. 抽象空间上的可加集合函数, Mat. Sbornik: I （第 1 章）, v. VIII (1940), pp.~307-348; II （第 2、3 章）, v. IX (1941), pp.~563-628; III （第 4 至 6 章）, v. XIII (1943), pp.~169-238.

4 P. Alexandroff-H. Hopf. Topologie I, Berlin (Springer), 1935.

5a Archimedis Opera quae quidem exstant omnia, nunc primus et gr. et lat. edita \ldots. Basilae, Jo. Hervagius, 1 卷，对开本, 1544.

5b Archimedis Opera Omnia. 3 vols., ed.~J. L. Heiberg, 2nd ed., Leipzig (Teubner), 1913-1915.

5c Les Oeuvres complètes d'Archimède. Transl. P. Ver Ecke, Paris-Brussels (Desclée-de Brouwer), 1921.

6 亚里士多德著作集，由 W. D. Ross 主持编译, Oxford, 1928 sqq.

7a E. Artin. 伽罗瓦理论 \ldots, Notre Dame, 1946.

7b E. Artin. Ueber die Zerlegung definiter Funktionen in Quadraten, Abh. Math. Sem. Univ. Hamburg, v. V (1927), pp.~100-115.

7c E. Artin. Zur Theorie der hypercomplexen Zahlen, Abh. Math. Sem. Univ. Hamburg, v. V (1927), pp.~251-260.

7d E. Artin. Einführung in die Theorie der Gammafunktion, Leipzig (Teubner), 1931.

8a E. Artin und O. Schreier. Algebraische Konstruktion reeler Körper, Abh. Math. Sem. Univ. Hamburg, v. V (1927), pp.~85-99.

8b E. Artin und O. Schreier. Eine Kennzeichnung der reell abgeschlossenen Körper, Abh. Math. Sem. Univ. Hamburg, v. V (1927), pp.~225-231.

9a C. Arzelà. Sulla integrabilità di una serie di funzioni, Rendic. Acc. dei Lincei (4), v. I (1885), pp.~321-326.

9b C. Arzelà. Sulla integrazione per serie. Rendic. Acc. dei Lincei (4), v. I (1885), pp.~532-537 与 566-569.

10 G. Ascoli. Le curve limiti di una varietà data di curve, Mem. Acc. dei Lincei (3), v. XVIII (1883), pp.~521-586.

11a R. Baire. Leçons sur les fonctions discontinues, Paris (Gauthier-Villars), 1905.

11b R. Baire. Sur les fonctions de variables réelles, Ann. di Mat. (3), v. III (1899), pp.~1-123.

12 R. Baire, E. Borel, J. Hadamard, H. Lebesgue. Cinq lettres sur la théorie des ensembles, Bull. Soc. Math. de France, v. XXXIII (1905), pp.~261-273.

13 H. F. Baker. 交错式与连续群, Proc. Lond. Math. Soc. (2), v. III (1905), pp.~24-47.

14a S. Banach. Théorie des opérations linéaires, Warszawa, 1932.

14b S. Banach. Sur le problème de la mesure, Fund. Math., v. IV (1923), pp.~7-33.

14c S. Banach. Sur les fonctionnelles linéaires, Studia Math., v. I (1929), pp.~211-216 与 223-239.

15 S. Banach 与 H. Steinhaus. Sur le principe de condensation des singularités, Fund. Math., v. IX (1927), pp.~50-61.

16a I. Barrow. Lectiones Geometricae \ldots, Londini, 1670.

16b I. Barrow. 数学著作集, ed.~W. Whewell, Cambridge, 1860.

17a O. Becker. Eudoxos-Studien, Quellen und Studien zur Geschichte der Math., Abt. B: Studien, v. II (1933), pp.~311-333 与 369-387，以及 v. III (1936), pp.~236-244 与 370-388.

17b O. Becker. Die Lehre vom Geraden und Ungeraden im Neunten Buch der Euklidischen Elementen, Quellen und Studien zur Geschichte der Math., Abt. B : Studien, v. III (1936), pp.~533-553.

18a E. Beltrami. Saggio di interpretazione della geometrica non-euclidean, Giorn. di Mat., v. VI (1868), pp.~284-312.

18b E. Beltrami. Teoria fondamentale degli spazii di curvatura costante, Ann. di Mat. (2), v. II (1868-69), pp.~232-255.

19a Jakob Bernoulli. Opera, 2 vol;., Geneva (Cramer-Philibert), 1744.

19b Jakob Bernoulli. Ars Conjectandi, Basel, 1713 (transl. S. Vastel, Paris, 1801).

20a Johann Bernoulli. Opera Omnia, 4 vol., Lausanne-Geneva (m. M. Bousquet), 1742.

20b Johann Bernoulli. Der erste Integralrechnung (Ostwald's Klassiker, no 194), Leipzig-Berlin (Engelmann), 1914.

20c Johann Bernoulli. Die Differentialrechnung (Ostwald's Klassiker, no 211), Leipzig (Akad. Verlag), 1924.

21 E. Bezout. Théorie générale des équations algébriques, Paris, 1779.

22a G. Birkhoff. 连续群与线性空间, Rec. Math. Moscow, v. I, (1936), pp.~635-642.

22b G. Birkhoff. 李代数与李群的矩阵可表示性, 《数学年鉴》, v. XXXVIII (1937), pp.~526-532.

23 J. M. Bochenski. 古代形式逻辑, 《逻辑研究丛书》, Amsterdam (North-Holland Publ. Co.), 1951.

24a S. Bochner. Integration von Funktionen deren Werte die Elemente eines Vektorraumes sind, Fund. Maht., v. XX (1933), 262-276.

24b S. Bochner. 调和分析与概率论, Berkeley (加州大学出版社), 1960.

25 P. Böhner. 中世纪逻辑，其自 1250 年至约 1400 年的发展概要, Chicago, 1952.

26 H. Bohr und J. Mollerup. Laerebog i matematisk Analyse, v. III, Kopenhagen, 1922.

27a B. Bolzano. Oeuvres, 5 vol., Prague, 1930-1948.

27b B. Bolzano. Paradoxien des Unendlichen, Leipzig, 1851 （英译本，New Haven (Yale Univ. Press), 1950）.

27c B. Bolzano. Rein Analytischer Beweis der Lehrsatzes, dass zwischen zwei Werthen, die ein entgegengesetzes Resultat gewahren, wenigstens eine reelle Wurzel liegt (Ostwald's Klassiker, no 153), Leipzig, 1905.

28a R. Bombelli. L'Algebra, Bologna (G. Rossi), 1579.

28b R. Bombelli, L'Algebra, libri IV e V, ed E. Bortolotti, Bologna (Zanichelli), 1929.

29 G. Boole. 逻辑著作集, 2 vol., ed.~P. Jourdain, Chicago- London, 1916.

30 A. Borel. Groupes linéaires algébriques, 《数学年鉴》, v. LXIV (1956), pp.~20-80.

31 A. Borel 与 J. Tits. Groupes réductifs, Publ. math. I. H. E. S., no 27 (1965), pp.~55-150.

32a E. Borel. Leçons sur la théorie des fonctions, Paris (Gauthier-Villars), 1898.

32b E. Borel. Les probabilités dénombrables et leurs applications arithmétiques, Rend. Circ. Mat. Palermo, v. XXVII (1909), pp.~286-310.

33 H. Bosmans. Sur le ``libro del Algebra'' de Pedro Nuñez, Bibl. Math. (3), v. VIII (1907-08), pp.~154-169.

34a R. Brauer. Über Systeme hypercomplexer Zahlen, Math. Zeitschr., v. XXX (1929), pp.~79-107.

34b R. Brauer. Eine Bedingung für vollständige Reduzibilität von Darstellungen gewöhnlicher und infinitesimaler Gruppen, Math. Zeitschr., v. XLI (1936), pp.~330-339.

35 A. Bravais. Mémoires sur les polyèdres de forme symétrique, Journ. de Math. (1), v. XIV (1849), pp.~141-180.

36 H. Briggs. Arithmetica logarithmica, London, 1624.

37 A. Brill und M. Noether. Ueber algebraische Funktionen, Math. Ann., v. VII (1874), pp.~269-310.

38 Lord Brouncker. 用有理数无穷级数求双曲线面积 \ldots, Phil. Trans., v. III (1668), pp.~645-649.

39a L. E. J. Brouwer. 直觉主义与形式主义, Bull. Amer. Math. Soc., v. XX (1913), pp.~81-96.

39b L. E. J. Brouwer. Zur Begründung des intuitionistischen Mathematik, Math. Ann., v. XCIII (1925), pp.~244-257, v. XCV (1926), pp.~453-473, v. XCVI (1926), pp.~451-458.

40 F. Bruhat. Représentations induites des groupes de Lie semi-simples complexes, C. R. Acad. Sci., v. CCXXXVIII (1954), pp.~437-439.

41 F. Bruhat 与 J. Tits. Groupes algébriques simples sur un corps local, Proc. Conf. on Local Fields, pp.~23-36, Berlin (Springer), 1967.

42 C. Burali-Forti. Sopra un theorema del Sig. G. Cantor, Atti Acad. Torino, v. XXXII (1896-97), pp.~229-237.

43 H. Burkhardt. Die Anfänge der Gruppentheorie und Paolo Ruffini, Zeitschr. für Math. und Phys., v. XXXVII (1892), Suppl., pp.~121-159.

44a W. Burnside. 论任意线性代换群的可约性条件, Proc. Lond. Math. Soc., v. III (1905), pp.~430-434.

44b W. Burnside. 有限阶群论, 2nd ed., Cambridge (University Press), 1911.

45 W. H. Bussey. 数学归纳法的起源. Amer. Math. Monthly, v. XXIV (1917), pp.~199-207.

46a J. E. Campbell. 论与连续变换群理论有关的算子合成规律, Proc. Lond. Math. Soc., (1), v. XXVIII (1897), pp.~381-390.

46b J. E. Campbell. 论算子的合成规律（第二篇）, Proc. Lond. Math. Soc., (1), v. XXIX (1898), pp.~14-32.

47 G. Cantor. Gesammelte Abhandlungen, Berlin (Springer), 1932.

48 G. Cantor, R. Dedekind. Briefwechsel, ed.~J. Cavaillès-E. Noether, Actual. Scient. et Ind., no 518, Paris (Hermann), 1937.

49 C. Caratheodory. Vorlesungen über reelle Funktionen, Leipzig-Berlin (Teubner), 1918.

50 H. Cardano. Opera, 10 vol., Lyon, 1663.

51 L. Carnot. Géométrie de Position, Paris, 1803.

52a E. Cartan. Oeuvres complètes, 6 vol.~(3 parts), Paris (Gauthier-Villars), 1953-55.

52b E. Cartan. Leçons sur la théorie des spineurs (Actual. Scient. et Ind., nos 643 与 701), Paris (Hermann), 1938.

53 H. Cartan. Théorie des filtres, C. R. Acad. Sci., v. CCV (1937), pp.~595-598; Filtres et ultrafiltres, ibid., pp.~777-779.

54 R. Carter. 单群与单李代数, Journ. Lond. Math. Soc., v. XL (1965), pp.~193-240.

55 H. Casimir und B. L. Van der Waerden. Algebraischer Beweis der vollständigen Reduzibilität der Darstellungen halbeinfacher Liescher Gruppen, Math. Ann., v. CXI (1935), pp.~1-12.

56a A.-L. Cauchy. Oeuvres complètes, 26 vol., （2 辑）, Paris (Gauthier-Villars), 1882-1958.

56b Leçons de calcul différentiel et de calcul intégral, rédigées principalement d'après les méthodes de M. A.-L. Cauchy，由 l'Abbé Moigno 编写，v. II, Paris, 1844.

57a B. Cavalieri. Geometria indivisibilibus continuorum quadam ratione promota, Bononiae, 1635 (2nd ed., 1653).

57b B. Cavalieri. Exercitationes geometricae sex, Bononiae, 1647.

58 A. Cayley. 数学论文集, 13 vol., Cambridge (University Press), 1889-1898.

59 L. Cesari. 曲面面积，Princeton, 1954.

60a M. Chasles. Note sur les propriétés générales du système de deux corps, Bull. de Férussac, v. XIV (1830), pp.~321-326.

60b M. Chasles. Aperçu historique sur l'origine et le développement des méthodes en géometrie, Brussels, 1837.

61 C. Chabauty. Sur les points rationnels des courbes algébriques de genre supérieur à l'unité, C. R. Acad. Sci., v. CCXII (1941), pp.~882-884.

62a C. Chevalley. 旋量的代数理论, New-York (Columbia University Press), 1954.

62b C. Chevalley. Généralisation de la théorie du corps de classes pour les extensions infinies, Journ. de Math., (9), v. XV (1936), pp.~359-371.

62c C. Chevalley. 论局部环理论, 《数学年鉴》, v. XLIV (1943), pp.~690-708.

62d C. Chevalley. 李群论, Princeton (University Press), 1946.

62e C. Chevalley. Sur la classification des algèbres de Lie simples et de leurs représentations, C. R. Acad. Sci., v. CCXXVII (1948), pp.~1136-1138.

62f C. Chevalley. 由反射生成的有限群的不变量, 《美国数学杂志》, v. LXXVII (1955), pp.~778-782.

62g C. Chevalley. Sur certains groupes simples, Tohōku Math. Journ., (2), v. VII (1955), pp.~14-66.

62h C. Chevalley. Classification des groupes de Lie algébriques, 2 vol., Paris (Inst. H. Poincaré), 1956-58.

63 G. Choquet. 容度理论, Ann. Inst. Fourier, v. V (1953-54), pp.~131-295.

64 N. Chuquet. Le Triparty en la Science des Nombres, ed.~A. Marre, Bull. bibl. storia math., v. XIII (1880), pp.~555-659 与 693-814.

65 W. K. Clifford. 数学论文，London (Macmillan), 1882.

66 I. Cohen 与 A. Seidenberg. 素理想与整相关性, Bull. Amer. Math. Soc., v. LII (1946), pp.~252-261.

67 P. J. Cohen. 连续统假设的独立性, Proc. Nat. Acad. Sci. U. S. A., v. L (1963), pp.~1143-1148 与 v. LI (1964), pp.~105-110.

68 A. J. Coleman. 单群的贝蒂数, 《加拿大数学杂志》, v. X (1958), pp.~349-356.

69a L. Couturat. La logique de Leibniz d'après des documents inédits, Paris (Alcan), 1901.

69b L. Couturat. La Philosophie des mathématiques de Kant, Revue de Métaph. et de Morale, v. XII (1904), pp.~321-383.

70a H. S. M. Coxeter. 基本区域为单形的群, Journ. Lond. Math. Soc., v. VI (1931), pp.~132-136.

70b H. S. M. Coxeter. 具有正棱柱形图形的多胞形, Proc. Lond. Math. Soc., (2), v. XXXIV (1932), pp.~126-189.

70c H. S. M. Coxeter. 由反射生成的离散群, 《数学年鉴》, v. XXXV, (1934), pp.~588-621.

70d H. S. M. Coxeter. 形如 \(R_{i}^{2} = (R_{i}, R_{j})^{k_{ij}} = 1\) 的有限群的完全枚举，Journ. Lond. Math. Soc., v. X (1935), pp.~21-25.

70e H. S. M. Coxeter 载于 H. Weyl. 连续群的结构与表示 （高等研究院， N. Jacobson 与 R. Brauer 记录的油印讲义, 1934-35）：附录。

70f H. S. M. Coxeter. 正多胞形, New York (Macmillan), 1948 (2nd ed., 1963).

70g H. S. M. Coxeter. 由反射生成的有限群的生成元之积, Duke Math. Journ., v. XVIII (1951), pp.~765-782.

71 H. S. M. Coxeter 与 W. O. J. Moser. 离散群的生成元与关系, Ergebn. der Math., Neue Folge, Bd. 14, Berlin (Springer), 1957 (2nd ed., 1965).

72 G. Cramer. Introduction à l'analyse des lignes courbes, Geneva (Cramer 与 Philibert), 1750.

73 H. Curry. 形式主义数学哲学概要, Amsterdam (North-Holland Publ. Co.), 1951.

74 M. Curtze. Über die Handschrift ``Algorismes proportionum magistri Nicolay Orem'', Zeitschr. für Math. und Phys., v. XIII, Supplem. (1868), pp.~65-79 与 pp.~101-104.

75a J. d'Alembert. Encyclopédie, Paris, 1751-1765.

75b J. d'Alembert. Sur les principes métaphysiques de Calcul infinitésimal, Mélanges de litt., d'hist. et de philosophie, new ed., v. V, Amsterdam (1768), pp.~207-219.

75c J. d'Alembert. Misc. Taur., v. III (1762-65), p.~381.

76a P. J. Daniell. 积分的一般形式, 《数学年鉴》 (2), v. XIX (1918), pp.~279-294.

76b P. J. Daniell. 无穷多个维数中的积分. 《数学年鉴》, (2), v. XX (1919), pp.~281-288.

76c P. J. Daniell. Stieltjes-Volterra products, Congr. Intern. des Math., Strasbourg, 1920, pp.~130-136.

76d P. J. Daniell. 无穷多个维数中的有界变差函数, 《数学年鉴》, (2), v. XXI (1919-20), pp.~30-38.

77 G. Darboux. Mémoire sur les fonctions discontinues, Ann. Ec. Norm. Sup. (2), v. IV (1875), pp.~57-112.

78 B. Datta 与 A. N. Singh. 印度数学史, 2 vol., Lahore (Motilal Banarsi Das), 1935-38.

79 R. Dedekind. Gesammelte mathematische Werke, 3 vol., Braunschweig (Vieweg), 1932.

80 R. Dedekind und H. Weber. Theorie der algebraischen Funktionen einer Veränderlichen, J. de Crelle, v. XCII (1882), pp.~181-290.

81 M. Dehn. Über den Rauminhalt, Math. Ann., v. LV (1902), pp.~465-478.

82 Ch. de la Vallée Poussin. Intégrales de Lebesgue, Fonctions d'ensembles, Classes de Baire, Paris (Gauthier-Villars), 1916 (2nd ed., 1936).

83a A. de Morgan. 论三段论（III）, Trans. Camb. Philos. Soc., v. X (1858), pp.~173-230.

83b A. de Morgan. 论三段论（IV）与关系逻辑, Trans. Camb. Philos. Soc., v. X (1860), pp.~331-358.

84 G. Desargues. Oeuvres, ed.~Pourra, v. I, Paris (Leiber), 1864.

85a R. Descartes. Oeuvres, ed.~Ch. Adam 与 P. Tannery, 13 vol., Paris (L. Cerf), 1897-1913.

85b R. Descartes. Geometria，Fr.~van Schooten 拉丁文译本，2nd ed., 2 vol., Amstardam (Elzevir), 1659-61.

86 J. A. de Séguier. Théorie des groupes finis. Éléments de la théorie des groupes abstraits, Paris (Gauthier-Villars), 1904.

87 M. Deuring. Algebren (Erg. der Math., Bd. 4), Berlin (Springer), 1937.

88a L. E. Dickson. 线性结合代数与阿贝尔方程, Trans. Amer. Math. Soc., v. XV (1914), pp.~31-46.

88b L. E. Dickson. 任意域上的线性群理论, Trans. Amer. Math. Soc., v. II (1901), pp.~363-394.

88c L. E. Dickson. 单群的一个新系统, Math. Ann., v. LX (1905), pp.~137-150.

88d L. E. Dickson. 任意域中与三次曲面上 27 条直线的构形相关的一类群, Quart. Journ. Pure 与 App. Math., v. XXXIII (1901), pp.~145-173 与 v. XXXIX (1908), pp.~205-209.

89 H. Diels. Die Fragmente der Vorsokratiker, 2te Aufl., 2 vol., Berlin (Weidmann), 1906-07.

90a J. Dieudonné. Une généralisation des espaces compacts, Journ. de Math., (9), v. XXIII (1944), pp.~65-76.

90b J. Dieudonné. La géométrie des groupes classiques (Erg. der Math. Neue Folge, Heft 5), Berlin-Göttingen-Heidelberg (Springer), 1955.

91a Diophanti Alexandrini Opera Omnia \ldots. 2 vol., ed.~P. Tannery, Lipsiae (Teubner), 1893-95.

91b Diophante d'Alexandrie. Transl. P. Ver Eecke, Bruges (Desclée de Brouwer), 1926.

92a P. G. Lejeune-Dirichlet. Werke, 2 vol., Berlin (Reimer), 1889-1897.

93 J. L. Doob. 函数空间中的概率, Bull. Amer. Math. Soc., v. LIII (1947), pp.~15-30.

94a P. du Bois-Reymond. Sur la grandeur relative des infinis des fonctions, Ann. di Mat. (2), v. IV (1871), pp.~338-353.

94b P. du Bois-Reymond. Üeber asymptotische Werthe, infinitäre Approximationen und infinitäre Auflösung von Gleichungen, Math. Ann., v. VIII (1875), pp.~362-414.

95 N. Dunford. 线性空间中的一致性, Trans. Amer. Math. Soc., v. XLIV (1938), pp.~305-356.

96 N. Dunford 与 B. Pettis. 可和函数上的线性运算, Trans. Amer. Math. Soc., v. XLVII (1940), pp.~323-392.

97 W. Dyck. Gruppentheoretische Studien, Math. Ann., v. XX (1882), pp.~1-44.

98a E. Dynkin. Campbell-Hausdorff 公式系数的计算（俄文）, Dokl. Akad. Nauk., v. LVII (1947), pp.~323-326.

98b E. Dynkin. 赋范李代数与解析群（俄文）, Uspehi Mat. Nauk, v. V (1950), pp.~135-186（英文译本：Amer. Math. Soc. Transl., (1) vol.~9, pp.~470-534）.

99 D. Egoroff. Sur les suites de fonctions mesurables, C. R. Acad. Sci., v. CLII (1911), pp.~244-246.

100 M. Eichler. Quadratische Formen und orthogonale Gruppen, Berlin-Göttingen-Heidelberg (Springer), 1952.

101 A. Einstein. 布朗运动理论研究, New York (Dover), 1956.

102a G. Eisenstein. Beweis der Reciprocitätsgesetze für die cubischen Reste in der Theorie der aus dritten Wurzeln der Einheit zusammengesetzen Zahlen, J. de Crelle, v. XXVII (1844), pp.~289-310.

102b G. Eisenstein. Zur Theorie der quadratischen Zerfällung der Primzahlen 8n + 3, 7n + 2 und 7n + 4, J. de Crelle, v. XXXVII (1848), pp.~97-126.

102c G. Eisenstein. Über einige allgemeine Eigenschaften der Gleichung von welcher die Teilung der ganzen Lemniscate abhängt, nebst Anwendung derselben auf die Zahlentheorie, J. de Crelle, v. XXXIX (1850), pp.~160-179 与 224-287.

103 G. Eneström. Kleine Bemerkung zur letzten Auflage von Cantors Vorlesungen zur Geschichte der Mathematik, Bibl. Math. (3), v. VIII (1907), pp.~412-413.

104a F. Engel. Über die Definitionsgleichung der continuierlichen Transformationsgruppen, Math. Ann. v. XXVII (1886), pp.~1-57.

104b F. Engel. Die Erzeugung der endlichen Transformationen einer projectiven Gruppe durch die infinitesimalen Transformationen der Gruppe, I, Leipziger Ber., v. XLIV (1892), pp.~279-296; II (mit Beiträgen von E. Study), ibid., v. XLV (1893), pp.~659-696.

105a F. Engel und P. Stäckel. Die Theorie der Parallellinien von Euklid bis auf Gauss, Leipzig (Teubner), 1895.

105b F. Engel und P. Stäckel. Urkunden zur Geschichte der nichteuklidischen Geometrien, 2 vol., Leipzig (Teubner), 1898-1913.

106 Enzyklopädie der Mathematischen Wissenschaften. 1te Aufl., 20 vol., Leipzig (Teubner), 1901-1935.

107 Euclidis Elementa. 5 vol., ed.~J. L. Heiberg, Lipsiae (Teubner), 1883-88.

108a L. Euler. Opera Omnia, 已出 46 卷（3 辑）, Leipzig-Berlin-Zürich (Teubner 与 O. Füssli), 1911-1957.

108b L. Euler. Formulae generales pro translatione quacunque corporum rigidorum, Novi Comm. Acad. Sc. imp. Petrop., v. XX (1776), pp.~189-207.

109 P. Fermat. Oeuvres, 5 vol., Paris (Gauthier-Villars), 1891-1922.

110 X. Fernique. Processus linéaires, processus généralisés, Ann. Inst. Fourier, v. XVII (1967), pp.~1-92.

111 E. Fischer. Sur la convergence en moyenne, C. R. Acad. Sci., v. CXLIV (1907), pp.~1022-1024.

112 J. B. Fourier. Oeuvres, 2 vol., Paris (Gauthier-Villars), 1888-90.

113a A. Fraenkel. Über die Teiler der Null und die Zerlegung von Ringen, J. de Crelle, v. CXLV (1914), pp.~139-176.

113b A. Fraenkel. Zu den Grundlagen der Cantor-Zermeloschen Mengenlehre, Math. Ann., v. LXXXVI (1922), pp.~230-237.

113c A. Fraenkel. Einleitung in die Mengenlehre, 3te Aufl., Berlin (Springer), 1928.

114 D. C. Fraser. 牛顿插值公式, Journ. Inst. Actuaries, v. LI (1918), pp.~77-106 与 pp.~211-232 与 v. LVIII (1927), pp.~53-95 （诸篇重印为一卷，London, s. d.）.

115a M. Fréchet. Sur quelques points du calcul fonctionnel, Rend. Circ. Mat. Palermo, v. XXII (1906), pp.~1-74.

115b M. Fréchet. Les ensembles abstraits et le Calcul fonctionnel, Rend. Circ. Mat. Palermo, v. XXX (1910), pp.~1-26

115c M. Fréchet. Sur l'intégrale d'une fonctionnelle étendue à un ensemble abstrait, Bull. Soc. Math. de France, v. XLIII (1915), pp.~248-265.

115d M. Fréchet. Les familles et fonctions additives d'ensembles abstraits, Fund. Math., v. IV (1923), pp.~229-265 与 v. V (1924), pp.~206-251.

116 I. Fredholm. Sur une classe d'équations fonctionnelles, Acta Mathematica, v. XXVII (1903), pp.~365-390.

117a G. Frege. Begriffsschrift eine der arithmetischen nachgebildete Formelsprache des reinen Denkens, Halle, 1879.

117b G. Frege. Die Grundlagen der Arithmetik, 第 2 版附 J. L. Austin 的英译, New-York, 1950.

117c G. Frege. Grundgesetze der Arithmetik, begriffschriftlich abgeleitet, 2 vol., Iena, 1893-1903.

118 R. Fricke und F. Klein. Theorie der automorphen Funktionen, Leipzig (Teubner), 1897.

119 G. Frobenius. Gesammelte Abhandlungen (ed.~J. P. Serre), 3 vol., Berlin-Heidelberg-New York (Springer), 1968.

120 G. Frobenius und L. Stickelberger. Ueber Gruppen von vertauschbaren Elementen, J. de Crelle, v. LXXXVI (1879), pp.~217-262.

121 G. Fubini. Sugli integrali multipli, Rendic. Acc. dei Lincei (5), v. XVI (1907), pp.~608-614.

122a Galileo Galilei. Discorsi e Dimostrazioni \ldots, Leyden (Elzevir), 1638.

122b Galileo Galilei. Opere, Ristampa della Ed. Nazionale, 20 vol., Firenze (Barbera), 1929-1939.

123 E. Galois. Écrits et mémoires mathématiques (ed.~R. Bourgne 与 J. Y. Azra), Paris (Gauthier-Villars), 1962.

124a C. F. Gauss. Werke, 12 vol., Göttingen, 1870-1927.

124b Die vier Gauss'schen Beweise für die Zerlegung ganzer algebraischer Functionen in reelle Factoren ersten oder zweiten Grades (Ostwald's Klassiker, no 14), Leipzig (Teubner), 1904.

125a I. Gelfand. Abstrakte Funktionen und lineare Operatoren, Mat. Sborn. (N.S.), v. IV (1938), pp.~235-284.

125b I. Gelfand. 广义随机过程（俄文）, Dokl. Akad. Nauk, v. C (1955), pp.~853-856.

125c I. Gelfand. 泛函分析的若干问题（俄文）, Uspehi Mat. Nauk, v. XI (1956), pp.~3-12（英文译本：Amer. Math. Soc. Transl., (2), vol.~16 (1960), pp.~315-324）.

126 I. Gelfand 与 N. Ya Vilenkin. 广义函数, vol.~IV, New York (Academic Press), 1964.

127 G. Gentzen. Die gegenwärtige Lage in der mathematischen Grundlagenforschung. Neue Fassung des Widerspruchsfreiheitsbeweises für die reine Zahlentheorie (Forschungen zur Logik \ldots, Heft 4, Leipzig (Hirzel), 1938).

128 Giorgini. Sopra alcune proprietà de' piani de' momenti \ldots, Mem. Soc. Ital. d.~Sc. res.（Modena）, v. XX (1828), pp.~243-254.

129 A. Girard. Invention nouvelle en Algèbre, Amsterdam, 1629 （重印，ed.~Bierens de Haan, Leyde, 1884）.

130a K. Gödel. Über formal unentscheidbare Sätze der Principia Mathematica und verwandter Systeme, Monatsh. für Math. und Phys., v. XXXVIII (1931), pp.~173-198.

130b K. Gödel. 选择公理与广义连续统假设的相容性 (《数学年鉴》研究丛书, no 3), Princeton, 1940.

130c K. Gödel. 什么是康托尔的连续统假设? Amer. Math. Monthly, v. LIV (1947), pp.~515-525.

131 D. Gorenstein. Finite groups, New York (Harper 与 Row), 1968.

132 E. Goursat. Sur les substitutions orthogonales et les divisions régulières de l'espace, Ann. Ec. Norm. Sup., (3), v. VI (1889), pp.~9-102.

133 J. P. Gram. Ueber die Entwicklung reeller Functionen in Reihen mittelst der Methode der kleinsten Quadrate, J. de Crelle, v. XCIV (1883), pp.~41-73.

134 H. Grassmann. Gesammelte Werke, 3 vol., Leipzig (Teubner), 1894-1911.

135 P. Gregorii a Sancto Vicentio. Opus Geometricum Quadraturae Circuli et Sectionum Coni \ldots, 2 vol., Antverpiae, 1647.

136a J. Gregory. Vera Circuli et Hyperbolae Quadraturae \ldots, Pataviae, 1667.

136b J. Gregory. Geometriae Pars Universalis, Pataviae, 1668.

136c J. Gregory. Exercitationes Geometricae, London, 1668.

136d James Gregory 三百周年纪念文集，收录他与 John Collins 的通信以及他此前未发表的数学手稿 \ldots. H. W. Turnbull 编, London (Bell 与 sons), 1939.

137 H. Grell. Beziehungen zwischen den Idealen verschiedener Ringe, Math. Ann., v. XCVII (1927), pp.~490-523.

138a A. Grothendieck. Éléments de géométrie algébrique, I, Publ. math. I. H. E. S., no 4 (1960).

138b A. Grothendieck. Produits tensoriels topologiques et espaces nucléaires, Mem. Amer. Math. Soc., no 16 (1955).

138c A. Grothendieck. Espaces vectoriels topologiques, 3rd ed., São Paulo (Publ. Soc. Mat. São Paulo), 1964.

139 A. Haar. Der Maassbegriff in der Theorie der kontinuierlichen Gruppen, 《数学年鉴》, v. XXXIV (1933), pp.~147-169.

140 J. Hachette 与 S. Poisson. Addition au mémoire précédent, Journ. de l'Ec. Polytechn., cahier 11 (an X), pp.~170-172.

141 J. Hadamard. Sur certaines applications possibles de la théorie des ensembles, Verhandl. Intern. Math. Kongress, Zürich, 1898, pp.~201-202.

142 H. Hahn. Ueber lineare Gleichungssysteme in linearen Räumen, J. de Crelle, v. CLVII (1927), pp.~214-229.

143 M. Hall. 自由李环以及自由群中高阶换位子的基, Proc. Amer. Math. Soc., v. I (1950), pp.~575-581.

144 P. Hall. 素数幂阶群理论的一项贡献, Proc. Lond. Math. Soc., (3), v. IV (1932), pp.~29-95.

145a W. R. Hamilton. 四元数讲义, Dublin, 1853.

145b W. R. Hamilton. 关于一种新单位根系的备忘录, Phil. Mag., (4), v. XII (1856), p.~446.

146 H. Hankel. Theorie der complexen Zahlensysteme, Leipzig (Voss), 1867.

147 G. H. Hardy. 无穷大的阶 (Cambridge tracts, no 12), 2nd ed., Cambridge University Press, 1924.

148 G. H. Hardy 与 E. M. Wright. 数论导引, Oxford, 1938.

149a Harish-Chandra. 论半单李代数的泛包络代数的若干应用, Trans. Amer. Math. Soc., v. LXX (1951), pp.~28-96.

149b Harish-Chandra. 论布吕阿的一个引理, Journ. de Math., (9), v. XXXV (1956), pp.~203-210.

150a H. Hasse. Ueber die Darstellbarkeit von Zahlen durch quadratischen Formen im Körper der rationalen Zahlen, J. de Crelle, v. CLII (1923), pp.~129-148.

150b H. Hasse. Ueber die Äquivalenz quadratischer Formen in Körper der rationalen Zahlen, J. de Crelle, v. CLII (1923), pp.~205-224.

150c H. Hasse. Kurt Hensels entscheidender Anstoss zur Entdeckung des Lokal-Global-Prinzips, J. de Crelle, v. CCIX (1960), pp.~3-4.

150d H. Hasse. Zahlentheorie, Berlin (Akad. Verlag), 1949.

151 H. Hasse und H. Scholz. Die Grundlagenkrise der grieschischen Mathematik, Charlottenburg (Pan-Verlag), 1928 (= Kant-Studien, v. XXXIII (1928), pp.~4-72).

152a F. Hausdorff. Die symbolische Exponentialformal in der Gruppentheorie, Leipziger Ber., v. LVIII (1906), pp.~19-48.

152b F. Hausdorff. Grundzüge der Mengenlehre, Leipzig (Veit), 1914.

152c F. Hausdorff. Mengenlehre, Berlin (de Gruyter), 1927.

153a T. Heath. 希腊数学史, 2 vol., Oxford 1921.

153b T. Heath. 佩尔加的阿波罗尼奥斯，圆锥曲线专论, Cambridge University Press, 1896.

153c T. Heath. 阿基米德的方法, Cambridge, 1912.

153d T. Heath. 亚里士多德中的数学, Oxford (Clarendon Press), 1949.

153e T. Heath. 欧几里得《几何原本》的十三卷 \ldots, 3 vol., Cambridge, 1908.

153f T. Heath. 亚历山大里亚的丢番图, 2nd ed., Cambridge, 1910.

154a E. Heine. Ueber trigonometrische Reihen, J. de Crelle, v. LXXI (1870), pp.~353-365.

154b E. Heine. Die Elemente der Functionenlehre, J. de Crelle, v. LXXIV (1872), pp.~172-188.

155 G. Heinrich. James Gregorys ``Vera circuli et hyperbolae quadratura'', Bibl. Math., (3), v. II (1901), pp.~77-85.

156 E. Helly. Ueber Systeme linearer Gleichungen mit unendlich vielen Unbekannten, Monatsh. für Math. und Phys., v. XXXI (1921), pp.~60-91.

157a K. Hensel. Über eine neue Begründung der Theorie der algebraischen Zahlen, Jahresber. der D. M. V., v. VI (1899), pp.~83-88.

157b K. Hensel. Ueber die Fundamentalgleichung und die ausserwesentlichen Diskriminantentheiler eines algebraischen Körpers, Gött. Nachr. (1897), pp.~254-260.

157c K. Hensel. Neue Grundlagen der Arithmetik, J. de Crelle, v. CXXVII (1902), pp.~51-84.

157d K. Hensel. Über die arithmetische Eigenschaften der algebraischen und transzendenten Zahlen, Jahresber. der D. M. V., v. XIV (1905), pp.~545-558.

157e K. Hensel. Ueber die arithmetischen Eigenschaften Zahlen, Jahresber. der D. M. V., v. XVI (1907), pp.~299-319, 388-393, 474-496.

157f K. Hensel. Theorie der algebraischen Zahlen, Leipzig (Teubner), 1908.

158 J. Herbrand. Recherches sur la théorie de la démonstration, Trav. Soc. Sci. Lett. Varsovie, cl. II (1930), pp.~33-160.

159 C. Hermite. Oeuvres, 4 vol., Paris (Gauthier-Villars), 1905-1917.

160 C. Hermite, T. Stieltjes. Correspondance, 2 vol., Paris (Gauthier-Villars), 1905.

161 J. Hessel. Krystallometrie oder Krystallonomie und Krystallographie （1830，重印于 Ostwald's Klassiker, nos 88 与 89, Leipzig (Teubner), 1897）.

162 A. Heyting. Mathematische Grundlagenforschung. Intuitionismus. Beweistheorie (Erg. der Math., Bd. 3), Berlin (Springer), 1934.

163a D. Hilbert. Gesammelte Abhandlungen, 3 vol., Berlin (Springer), 1932-35.

163b D. Hilbert. Grundzüge einer allgemeinen Theorie der Integralgleichungen, 2nd ed., Leipzig-Berlin (Teubner), 1924.

163c D. Hilbert. Grundlagen der Geometrie, 7th ed., Leipzig-Berlin (Teubner), 1930.

164 D. Hilbert und W. Ackermann. Grundzüge der theoretischen Logik, 3te Aufl., Berlin (Springer), 1949.

165 O. Hölder. Zurückführung einer beliebigen algebraischen Gleichung auf eine Kette von Gleichungen, Math. Ann., v. XXXIV (1889), pp.~26-56.

166 R. Hooke. 线性 \(p\) 进群及其李代数, 《数学年鉴》, v. XLIII (1942), pp.~641-655.

167 C. Hopkins. 具有左理想极小条件的环, 《数学年鉴》, (2), v. XL (1939), pp.~712-730.

168 A. Hurwitz. Über die Erzeugung der Invarianten durch Integration, Gött. Nachr., 1897, pp.~71-90 (= Math. Werke, v. II, pp.~546-564).

169a Christiani Hugenii, Zulichemii Philosophi vere magni, Dum viveret Zelemii Toparchae, Opera \ldots. 4 卷合订为 1 卷, Lugd. Batav., 1751.

169b C. Huygens. Oeuvres complètes, 22 vol., 海牙 (M. Nijhoff), 1888-1950.

170 N. Iwahori 与 H. Matsumoto. 论某些布吕阿分解与 \(p\) 进谢瓦莱群的海克环的结构, Publ. math. I. H. E. S., no 25 (1965), pp.~5-48.

171 C. G. J. Jacobi. Gesammelte Werke, 7 vol., Berlin (G. Reimer), 1881-91.

172a N. Jacobson. 李代数理论中的有理方法, 《数学年鉴》, v. XXXVI (1935), pp.~875-881.

172b N. Jacobson. 特征 p 的限制李代数的类，II, Duke Math. Journ., v. X (1943), pp.~107-121.

172c N. Jacobson. 任意环中本原理想集合的一个拓扑, Proc. Nat. Acad. Sci. U. S. A., v. XXXI (1945), pp.~333-338.

172d N. Jacobson. 环的结构 (Amer. Math. Soc. Coll. Public., v. 37), Providence, 1956.

173 B. Jessen. 无穷维空间中的积分理论, Acta Math., v. LXIII (1934), pp.~249-323.

174a C. Jordan. Traité des substitutions et des equations algébriques, 2nd ed., Paris (Gauthier-Villars et A. Blanchard), 1957.

174b C. Jordan. Mémoire sur les groupes de mouvements, Ann. di Mat., (2), v. II (1868-69), pp.~167-215 与 322-345 (= Oeuvres, v. IV, pp.~231-302, Paris (Gauthier-Villars 4), 1964.)

174c C. Jordan. Cours d'Analyse de l'École Polytechnique, 3rd ed., 3 vol., Paris (Gauthier-Villars), 1909-15.

175 H. Jung. Algebraischen Flächen, Hannover (Helwing), 1925.

176 J. P. Kahane. Séries de Fourier aléatoires, Sém. Bourbaki, no. 200, 第 12 年度, 1959-60, New York (Benjamin).

177 S. Kakutani. 关于无穷乘积测度空间的注记, II, Proc. Imp. Acad. Japan, v. XIX (1943), pp.~184-188.

178 I. Kant. Werke, ed.~E. Cassirer, 11 vol., Berlin (B. Cassirer), 1912-1921.

179a J. Kepler. Stereometria Doliorum, 1615.

179b J. Kepler. Neue Stereometrie der Fäser (Ostwald's Klassiker, no. 165), Leipzig (Engelmann), 1908.

180 W. Killing. Die Zusammensetzung der stetigen endlichen Transformationsgruppen : I) Math. Ann., v. XXXI (1888), pp.~252-290; II) ibid., v. XXXIII (1889), pp.~1-48; III) ibid., v. XXXIV (1889), pp.~57-122; IV) ibid., v. XXXVI (1890), pp.~161-189.

181 S. Kleene. 元数学导论, New York, 1952.

182 F. Klein. Gesammelte mathematische Abhandlungen, 3 vol., Berlin (Springer), 1921-23.

183 A. Kolmogoroff. Grundgebriffe der Wahrscheinlichkeitsrechnung (Erg. der Math., Bd. 2), Berlin (Springer), 1933.

184 G. Köthe. Neubegründung der Theorie der vollkommenen Räume, Math. Nachr., v. IV (1951), pp.~70-80.

185 E. Kötter. Die Entwicklung der synthetischen Geometrie, Leipzig (Teubner), 1901 (= Jahresbericht der D. M. V., v. V, 2tes Heft).

186a L. Kronecker. Werke, 5 vol., Leipzig (Teubner), 1895-1930.

186b L. Kronecker. Vorlesungen über die Theorie der Determinanten \ldots, Leipzig (Teubner), 1903.

187a W. Krull. Über verallgemeinerte endliche Abelsche Gruppen, Math. Zeitschr., v. XXIII (1925), pp.~161-196.

187b W. Krull. Theorie und Anwendung der verallgemeinerten Abelschen Gruppen, Sitzungsber. Heidelberger Akad. Wiss., 1926, no. 1, 32 pp.

187c W. Krull. Zur Theorie der allgemeinen Zahlringe, Math. Ann., v. XCIX (1928), pp.~51-70.

187d W. Krull. Galoische Theorie der unendlichen algebraischen Erweiterungen, Math. Ann., v. C (1928), pp.~687-698.

187e W. Krull. Primidealketten in allgemeine Ringbereichen, Sitzungsber. Heidelberg Akad. Wiss., 1928.

187f W. Krull. Allgemeine Bewertungstheorie, J. de Crelle, v. CLXVII (1931), pp.~160-196.

187g W. Krull. Beiträge zur Arithmetik kommutativer Integritätsbereiche III, Math. Zeitschr., v. XLII (1937), pp.~745-766.

187h W. Krull. Dimensionstheorie in Stellenringe, J. de Crelle, v. CLXXIX (1938), pp.~204-226.

187i W. Krull. Idealtheorie (Erg. der Math., Bd 4), Berlin (Springer), 1935.

188a E. Kummer. Sur les nombres complexes qui sont formés avec les nombres entiers réels et les racines de l'unité, Journ. de Math., (1), v. XII (1847), pp.~185-212.

188b E. Kummer. Zur Theorie der complexen Zahlen, J. de Crelle, v. XXXV (1847), pp.~319-326.

188c E. Kummer. Ueber die Zerlegung der aus Wurzeln der Einheit gebildeten complexen Zahlen in Primfactoren, J. de Crelle, v. XXXV (1847), pp.~327-367.

188d E. Kummer. Mémoire sur les nombres complexes composés de racines de l'unité et des nombres entiers, Journ. de Math., (1), v. XVI (1851), pp.~377-498.

188e E. Kummer. Über die allgemeinen Reciprocitätsgesetze unter den Resten und Nichtresten der Potenzen deren Grad eine Primzahl ist, Abhandl. der Kön. Akad. der Wiss. zu Berlin (1859), Math. Abhandl., pp.~19-159.

189a K. Kuratowski. Une méthode d'élimination des nombres transfinis des raisonnements mathématiques, Fund. Math., v. V (1922), pp.~76-108.

189b K. Kuratowski. Topologie, I, 2nd ed., Warszawa-Vroclaw, 1948.

190 J. Kürschak. Über Limesbildung und allgemeine Körpertheorie, J. de Crelle, v. CXLII (1913), pp.~211-253.

191 J. L. Lagrange. Oeuvres, 14 vol., Paris (Gauthier-Villars), 1867-1892.

192 E. Laguerre. Oeuvres, 2 vol., Paris (Gauthier-Villars), 1898-1905.

193 P. S. Laplace. Oeuvres, 14 vol., Paris (Gauthier-Villars), 1878-1912.

194 E. Lasker. Zur Theorie der Moduln und Ideale, Math. Ann., v. LX (1905), pp.~20-116.

195a M. Lazard. Quelques calculs concernant la formule de Hausdorff, Bull. Soc. Math. de France, v. XCI (1963), pp.~435-451.

195b M. Lazard. Groupes analytiques p-adiques, Publ. Math. I. H. E. S., no. 26 (1965), pp.~389-603.

196a H. Lebesgue. Intégrale, longueur, aire, Ann. di Mat. (3), v. VII (1902), pp.~231-259.

196b H. Lebesgue. Sur les séries trigonométriques, Ann. Ec. Norm. Sup., (3), v. XX (1903), pp.~453-485.

196c H. Lebesgue. Leçons sur l'Intégration et la recherche des fonctions primitives, Paris (Gauthier-Villars), 1904.

196d H. Lebesgue. Sur le problème de Dirichlet, Rend. Circ. Mat. Palermo, v. XXIV (1907), pp.~371-402.

196e H. Lebesgue. Sur l'intégration des fonctions discontinues, Ann. Ec. Norm. Sup. (3), v. XXVII (1910), pp.~361-450.

197 L. Le Cam. 随机过程的依分布收敛, Univ. Calif. Publ. Statistics, no. 11 (1957), pp.~207-236.

198a G. W. Leibniz. Mathematische Schriften, 7 vol., ed.~C. I. Gerhardt, Berlin-Halle (Ascher-Schmidt), 1849-63.

198b G. W. Leibniz. Philosophische Schriften, 7 vol., ed.~C. I. Gerhardt, Berlin, 1840-90.

198c G. W. Leibniz. Opuscules et fragments inédits, ed.~L. Couturat, Paris (Alcan), 1903.

198d Der Briefwechsel von Gottfried Wilhelm Leibniz mit Mathematikern. V. I, herausgegeben von C. I. Gerhardt, Berlin (Mayer und Müller), 1899.

199 B. Levi. Intorno alla teoria degli aggregati, R. Ist. Lombardo Sci. Lett. Rendic. (2), v. XXXV (1902), pp.~863-868.

200 E. E. Levi. Sulla struttura dei Gruppi finiti e continui, Atti Accad. Sci. Torino, v. XL (1905), pp.~551-565 (= Opere, v. I, p.~101-115).

201a P. Lévy. Leçons d'Analyse fonctionnelle, Paris (Gauthier-Villars), 1922) （第 2 版，题为：Problèmes concrets d'Analyse fonctionnelle，1951）.

201b P. Lévy. Processus stochastiques et mouvement brownien, Paris (Gauthier-Villars), 1948.

201c P. Lévy. Le mouvement brownien, Mémor. des Sci. Math., v. CXXVI (1954), Paris (Gauthier-Villars).

202 S. Lie. Gesammelte Abhandlungen, 7 vol., Leipzig (Teubner).

203 S. Lie und F. Engel. Theorie der Transformationsgruppen, 3 vol., Leipzig (Teubner), 1888-1893.

203bis J. Lindenstrauss und L. Tzafriri. 经典 Banach 空间, v. I, Berlin-Heidelberg-New York (Springer), 1977.

204a J. Liouville. Sur le développement des fonctions ou parties de fonctions en séries dont les divers termes sont assujettis à satisfaire à une même équation différentielle du second ordre contenant un paramètre variable, Journ. de Math. (1), v. I (1836), pp.~253-265 与 v. II (1837), pp.~16-35 与 418-436.

204b J. Liouville. D'un théorème dû à M. Sturm et relatif à une classe de fonctions transcendants, Journ. de Math. (1), v. I (1836), pp.~269-277.

204c J. Liouville. Sur des classes très étendues de quantités dont la valeur n'est ni algébrique ni même réductible à des irrationnelles algébriques, Journ. de Math. (1), v. XVI (1851), pp.~133-142.

205a R. Lipschitz. De explicatione per serias trigonometricas instituenda functionum unius variabilis arbitrariarum, J. de Crelle, v. LXIII (1864), pp.~296-308 （法文译者 P. Montel Acta Math., v. XXXVI (1912), pp.~261-295).

205b R. Lipschitz. Untersuchungen ueber die Summen von Quadraten, Bonn, 1886 （以法文摘要刊于 Bull. Sci. Math., v. XXI (1886), pp.~163-183).

206 N. Lobatschevsky. Pangeometrie (Ostwald's Klassiker, no. 130), Leipzig (Engelmann), 1902.

207 L. H. Loomis. 抽象调和分析导论, London-New York-Toronto (van Nostrand), 1953.

208 G. Loria. Le ricerche inedite de Evangelista Torricelli sopra la curve logaritmica, Bibl. Math. (3), v. I (1900), pp.~75-89.

209 N. Lusin. Leçons sur les ensembles analytiques et leurs applications, Paris (Gauthier-Villars), 1930.

210 E. Lutz. Sur l'équation \(y^{2} = x^{3} - Ax - B\) dans les corps \(p\) -adiques, J. de Crelle, v. CLXXVII (1937), pp.~237-247.

211 F. S. Macaulay. 论把给定模系分解为准素系，兼论希尔伯特数的若干性质, Math. Ann., v. LXXIV (1913), pp.~66-121.

212 A. Macbeath 与 S. Swierczkowski. 紧生成群中格的极限, 《加拿大数学杂志》, v. XII (1960), pp.~427-437.

213a G. W. Mackey. 论无穷维空间, Trans. Amer. Math. Soc., v. LVII (1945), pp.~155-207.

213b G. W. Mackey. 论凸拓扑空间, Trans. Amer. Math. Soc., v. LX (1946), pp.~519-537.

214 C. Maclaurin. 流数论, Edinburgh, 2 vol., 1742.

215a W. Magnus. Beziehungen zwischen Gruppen und Idealen in einen speziellen Ring, Math. Ann., v. CXI (1935), pp.~259-280.

215b W. Magnus. Über Beziehungen zwischen höheren Kommutatoren, J. de Crelle, v. CLXXVII (1937), pp.~105-115.

216 W. Magnus、A. Karass 与 D. Solitar. 组合群论, New York (Interscience), 1966.

217 D. Mahnke. Neue Einblicke in die Entdeckungsgeschichte der höheren Analysis, Abh. Preuss. Akad. der Wiss. Phys.-Math. Klasse, 1925, no. 1, 64 pp., Berlin, 1926.

218 F. Masères. Scriptores Logaritmici, 6 vol., London, 1791-1807.

219 L. Maurer. Über allgemeinere Invarianten-Systeme, Sitzungsber. München, v. XVIII (1888), pp.~103-150.

220 N. Mercator. Logarithmotechnia\ldots{} cui nunc accedit vera quadratura hyperbolae \ldots, Londini, 1668.

221a H. Minkowski. Gesammelte Abhandlungen, 2 vol., Leipzig-Berlin (Teubner), 1911.

221b H. Minkowski. Geometrie der Zahlen, Leipzig (Teubner), 1896.

222 R. A. Minlos. 广义随机过程及其测度扩张（俄文）, Trudy Mosk. Mat. Obschtsch., v. VIII (1959), pp.~497-518 (=《数学、统计学与概率译文选》, v. III (19), pp.~291-313, Amer. Math. Soc.).

223 A. F. Möbius. Gesammelte Werke, 4 vol., Leipzig (Hirzel), 1885-87.

224a T. Molien. Ueber Systeme höherer complexer Zahlen, Math. Ann., v. XLI (1893), pp.~83-156.

224b T. Molien. Über die Invarianten der linearen Substitutionsgruppen, Berliner Sitzungsber., 1897, pp.~1152-1156.

225 G. Monge. Géométrie descriptive, Paris, 1798 (new ed.~Paris (Gauthier-Villars), 1922).

226 D. Montgomery 与 L. Zippin. 拓扑变换群, New York (Interscience), 1955.

227 E. H. Moore 与 H. L. Smith. 极限的一般理论, 《美国数学杂志》, v. XLIV (1922), pp.~102-121.

228 S. G. Morley. 古代玛雅, Stanford University Press, 1946.

229 E. Nelson. 布朗运动的动力学理论, Mathematical Notes, Princeton (University Press), 1967.

230a J. Napier. Mirifici logarithmorum canonis descriptio, Lyon, 1619 （重印于 {[}218{]}, v. VI, pp.~475-623）.

230b Napier 三百周年纪念文集，London, 1915.

231 E. Netto. Zur Theorie der Elimination, Acta Math., v. VII (1885), pp.~101-104.

232 O. Neugebauer. Vorlesungen über die Geschichte der antiken Mathematik, Bd. I: Vorgrieschische Mathematik, Berlin (Springer), 1934.

233a I. Newton. Opuscula, 3 vol., Lausanne-Genève (M. Bousquet), 1744.

233b I. Newton. Philosophiae Naturalis Principia Mathematica, London, 1687 (new ed., Glasgow, 1871).

233c I. Newton. 自然哲学的数学原理, A. Motte 于 1729 年译成英文, 加州大学, 1946.

233d Isaac Newton 数学论文集. Vol. I: 1664-1666 (ed.~D. Whiteside), Cambridge (Univ. Press), 1967.

234 J. Nielsen. Die Isomorphismengruppen der freien Gruppen, Math. Ann., v. XCI (1924), pp.~169-209.

235 O. Nikodym. Sur une généralisation des intégrales de M. Radon, Fund. Math., v. XV (1930), pp.~131-179.

236a E. Noether. Idealtheorie in Ringbereichen, Math. Ann., v. LXXXIII (1921), pp.~24-66.

236b E. Noether. Abstrakter Aufbau der Idealtheorie in algebraischen Zahl- und Funktionenkörper, Math. Ann., v. XCVI (1926), pp.~26-61.

236c E. Noether. Hyperkomplexe Grössen und Darstellungstheorie, Math. Zeitschr., v. XXX (1929), pp.~641-692.

236d E. Noether. Nichtkommutative Algebra, Math. Zeitschr., v. XXXVII (1933), pp.~514-541.

237 E. Noether und R. Brauer. Über minimale Zerfallungskörper irreduzibler Darstellungen, Berliner Sitzungsber., 1927, pp.~221-228.

238 E. Noether und W. Schmeidler. Moduln in nichtkommutativen Bereichen, Math. Zeitschr., v. VIII (1920), pp.~1-35.

239 M. Noether. Über einen Satz aus der Theorie der algebraischen Funktionen, Math. Ann., v. VI (1873), pp.~351-359.

240 W. Osgood. 非一致收敛与级数的逐项积分, 《美国数学杂志》, v. XIX (1897), pp.~155-190.

241 P. Osmond. Isaac Barrow：其生平与时代, London, 1944.

242 A. Ostrowski. Über einige Lösungen der Funktionalgleichung \(\phi(x)\phi(y) = \phi(x.y)\) , Acta Math., v. XLI (1917), pp.~271-284.

243 R. E. A. C. Paley 与 N. Wiener. 复域中的傅里叶变换, Amer. Math. Soc. Coll. Publ., no. 19, New York, 1934.

244 B. Pascal. Oeuvres, 14 vol., ed.~Brunschvicg, Paris (Hachette), 1904-14.

245 M. Pasch und M. Dehn. Vorlesungen über neuere Geometrie, 2te Aufl., Berlin (Springer), 1926.

246a G. Peano. Applicazioni geometriche del calcolo infinitesimale, Torino, 1887.

246b G. Peano. Calcolo geometrico secondo l'Ausdehnungslehre di Grassmann, preceduto dalle operazioni della logica deduttiva, Torino, 1888.

246c G. Peano. Arithmeticas principia, novo methodo exposita, Torino, 1889.

246d G. Peano. I principii di Geometria, logicamente expositi, Torino, 1889.

246e G. Peano. Démonstration de l'intégrabilité des équations différentielles ordinaires, Math. Ann., v. XXXVII (1890), pp.~182-228.

246f G. Peano. Formulaire de Mathématiques, 5 vol., Torino, 1895-1905.

247 B. Peirce. 线性结合代数，《美国数学杂志》, v. IV (1881), pp.~97-221.

248a C. S. Peirce. 论数学的逻辑, 《美国艺术与科学学院院刊》, v. VII (1865-68), pp.~402-412.

248b C. S. Peirce. 论逻辑的代数, 《美国数学杂志》, v. III (1880), pp.~49-57.

248c C. S. Peirce. 论这些代数的相对形式, 《美国数学杂志》, v. IV (1881), pp.~221-225.

248d C. S. Peirce. 论除法无歧义的那些代数, 《美国数学杂志》, v. IV (1881), pp.~225-229.

248e C. S. Peirce. 论逻辑的代数, 《美国数学杂志》, v. VII (1884), pp.~190-202.

249 J. Pfanzagl 与 W. Pierlo. 紧集系, 《数学讲义》, no. 16 (1966), Berlin (Springer).

250 Plato. La République, transl. E. Chambry, 2 vol., Paris (Les Belles Lettres), 1932-49.

251a H. Poincaré. Oeuvres, 11 vol., Paris (Gauthier-Villars), 1916-1956.

251b H. Poincaré. Les méthodes nouvelles de la mécanique céleste, 3 vol., Paris (Gauthier-Villars), 1893-1899.

251c H. Poincaré. Science et hypothèse, Paris (Flammarion), 1902.

251d H. Poincaré. La valeur de la Science, Paris (Flammarion), 1905.

251e H. Poincaré. Science et méthode, Paris (Flammarion), 1908.

252 J. V. Poncelet. Traité des propriétés projectives des figures, 2 vol., 2nd ed., Paris (Gauthier-Villars), 1865.

253 L. S. Pontrjagin. 拓扑群，Princeton (Univ. Press), 1939.

254 Ju. V. Prokhorov. 随机过程的收敛与概率论中的极限定理, Theor. Probab. Appl., v. I (1956), pp.~156-214.

255 Ptolemaei Cl. Opera. Ed. J. L. Heiberg, 2 vol., Lipsiae (Teubner), 1898-1903 (transl. Halma, 重印, 2 vol., Paris (Hermann), 1927).

256 J. Radon. Theorie und Anwendungen der absolut additiven Mengenfunctionen, Sitzungsber. der math. naturwiss. Klasse der Akad. der Wiss. (Wien), v. CXXII, Abt. II a (1913), pp.~1295-1438.

257 G. Ricci et T. Levi-Civita. Méthodes de calcul différentiel absolu et leurs applications. Math. Ann., v. LIV (1901), pp.~125-201.

258 J. Richard. Les principes des Mathématiques et le problème des ensembles, Riv. Gén. des Sci. pures et appl., v. XVI (1905), pp.~541-543.

259a B. Riemann. Gesammelte mathematische Werke, 2nd ed., Leipzig (Teubner), 1892.

259b B. Riemann. Gesammelte Werke, Nachträge, Leipzig (Teubner), 1902.

259c B. Riemann. 载于 Lettere di E. Betti a P. Tardy, Rend. Accad. dei Lincei (5), v. XXIV \(^{1}\) (1915), pp.~517-519.

260a F. Riesz. Sur les systèmes orthogonaux de fonctions, C. R. Acad. Sci., v. CXLIV (1907), pp.~615-619.

260b F. Riesz. Stetigkeitsbegriff und abstrakte Mengenlehre, Atti del IV Congresso Intern. dei Matem., Roma, 1908, v. II, pp.~18-24.

260c F. Riesz. Untersuchungen über Systeme integrierbarer Funktionen, Math. Ann., v. LXIX (1910), pp.~449-497.

260d F. Riesz. Sur certains systèmes singuliers d'équations intégrales, Ann. Ec. Norm. Sup. (3), v. XXVIII (1911), pp.~33-62.

260e F. Riesz. Les systèmes d'équations linéaires à une infinité d'inconnues, Paris (Gauthier-Villars), 1913.

260f F. Riesz. Ueber lineare Funktionalgleichungen, Acta Math., v. XLI (1918), pp.~71-98.

260g F. Riesz. Zur Theorie des Hilbertschen Raumes, Acta litt. ac. scient. (Szeged), v. VII (1934-35), pp.~34-38.

260h F. Riesz. Sur quelques notions fondamentales dans la théorie générale des opérations linéaires, 《数学年鉴》 (2), v. XLI (1940), pp.~174-206.

261 M. Riesz. Sur le problème des moments, 3, Ark. für Math., v. XVII (1922-23), no. 16, 52 pp.

262 S. P. Rigaud. 科学人士通信集 \ldots, 2 vol., Oxford, 1841-42.

263 G. de Roberval. Ouvrages de Mathématique (Mémoires de l'Académie Royale des Sciences, v. VI, Paris (1730), pp.~1-478).

264 R. Robinson. 柏拉图对谬误的意识, Mind, v. LI (1942), pp.~97-114.

265 P. Ruffini. Opere Matematiche, 3 vol., Ed. Cremonese (Roma), 1953-54.

266 B. Russell 与 A. N. Whitehead. Principia Mathematica, 3 vol., Cambridge, 1910-13.

267 A. Rustow. Der Lügner, Diss. Erlangen, 1910.

268 S. Saks. 积分论, 2nd ed., New York (Stechert), 1937.

269 P. Alfonso Antonio de Sarasa. Solutio problematis \ldots, Antverpiae, 1649.

270 P. Samuel. La notion de multiplicité en Algèbre et en Géométrie algébrique, Journ. de Math. (9), v. XXX (1951), pp.~159-274.

271 G. Scheffers. Zurückführung complexer Zahlensysteme auf typische Formen, Math. Ann., v. XXXIX (1891), pp.~293-390.

272 E. Schering. Gesammelte mathematische Werke, 2 vol., Berlin (Mayer und Müller), 1902-1909.

273 L. Schäfli. Gesammelte mathematische Abhandlungen, 3 vol., Basel (Birkhäuser), 1950-56.

274a E. Schmidt. Zur Theorie der linearen und nichtlinearen Integralgleichungen. I. Teil: Entwicklung willkürlicher Funktionen nach Systeme vorgeschriebener, Math. Ann., v. LXIII (1907), pp.~433-476.

274b E. Schmidt. Ueber die Auflösung linearer Gleichungen mit unendlich vielen Unbekannten, Rend. Palermo, v. XXV (1908), pp.~53-77.

275a A. Schoenfleis. Entwicklung der Mengenlehre und ihrer Anwendungen, 2nd ed., Leipzig-Berlin (Teubner), 1913.

275b A. Schoenfleis. Die Krisis in Cantor's mathematischem Schaffen, Acta Math., v. L (1927), pp.~1-23.

276a O. Schreier. Abstrakte kontinuierliche Gruppen, Abh. math. Sem. Univ. Hambrug, v. IV (1926), pp.~15-32.

276b O. Schreier. Die Verwandschaft stetiger Gruppen in grossen, Abh. math. Sem. Univ. Hamburg, v. V (1927), pp.~233-244.

277 E. Schröder. Vorlesungen über die Algebra der Logik, 3 vol., Leipzig (Teubner), 1890.

278a F. Schur. Zur Theorie der aus Haupteinheiten gebildeten Komplexen, Math. Ann., v. XXXIII (1889), pp.~49-60.

278b F. Schur. Neue Begründung der Theorie der endlichen Transformationsgruppen, Math. Ann., v. XXXV (1890), pp.~161-197.

278c F. Schur. Zur Theorie der endlichen Transformationsgruppen, Math. Ann., v. XXXVIII (1891), pp.~273-286.

278d F. Schur. Über den analytischen Character der eine endliche continuierliche Transformationsgruppen darstellende Funktionen, Math. Ann., v. XLI (1893), pp.~509-538.

279a I. Schur. Über eine Klasse von Matrices, die sich einer gegebenen Matrix zuordnen lassen, Diss. Berlin, 1901.

279b I. Schur. Über die Darstellung der endlichen Gruppen durch gebrochene lineare Substitutionen, J. de Crelle, v. CXXVII (1904), pp.~20-50.

279c I. Schur. Neue Begründung der Theorie der Gruppencharaktere, Berliner Sitzungsber., 1905, pp.~406-432.

279d I. Schur. Arithmetische Untersuchungen über endliche Gruppen linearer Substitution, Berliner Sitzungsber., 1906, pp.~164-184.

279e I. Schur. Neue Anwendungen der Integralrechnung auf die Probleme der Invariantentheorie, Berliner Sitzungsber., 1924, pp.~189-208, 297-321, 346-355.

280 L. Schwartz. Théorie des distributions (Actual. Scient. et Ind., nos 1091 与 1122), Paris (Hermann), 1950-51.

281 I. Segal. 希尔伯特空间中的分布与典型算子系, Trans. Amer. Math. Soc., v. LXXXVIII (1958), pp.~12-41.

282 E. Selling. Ueber die idealen Primfactoren der complexen Zahlen, welche aus der Wurzeln einer beliebigen irreductiblen Gleichung rational gebildet sind, Zeitschr. für Math. und Phys., v. X (1865), pp.~17-47.

283a J.-P. Serre. Faisceaux algébriques cohérents, 《数学年鉴》, v. LXI (1955), pp.~197-278.

283b J.-P. Serre. Géométrie algébrique et géométrie analytique, Ann. Inst. Fourier, v. VI (1956), pp.~1-42.

284 J. A. Serret. Cours d'Algèbre Supérieure, 3rd ed., Paris (Gauthier-Villars), 1866.

285 C. L. Siegel. 辛几何, 《美国数学杂志》, v. LXV (1943), pp.~1-86.

286a T. Skolem. Einige Bemerkungen zur axiomatischen Begründung der Mengenlehre, Wiss. Vorträge, 5. Kongress Skand. Math., Helsingfors, 1922, pp.~217-232.

286b T. Skolem. Zur Theorie der associativen Zahlensysteme, Skr. norske Vid. Akad., Oslo, 1927, no. 12, 50 pp.

287 H. J. Smith. 数学论文集, 2 vol., Oxford, 1894.

288 Sommer. Introduction à la théorie des nombres algébriques (transl. A. Lévy), Paris (Hermann), 1911.

289 M. Souslin. Sur une définition des ensembles mesurables B sans nombres transfinis, C. R. Acad. Sci., v. CLXIV (1917), pp.~88-91.

290 E. Sparre-Andersen 与 B. Jessen. 论无穷乘积集合中测度的引入, Dansk. Vid. Selbskab. Mat. Fys. Medd., v. XXV (1948), no. 4, pp.~1-7.

291 A. Speiser. Theorie der Gruppen von endlicher Ordnung, 4th ed., Basel (Birkhäuser), 1956.

292 R. Steinberg. 有限反射群, Trans. Amer. Math. Soc., v. XCI (1959), pp.~493-504.

293 H. Steinhaus. Les probabilités dénombrables et leur rapport à la théorie de la mesure, Fund. Math., v. IV (1923), pp.~286-310.

294a E. Steinitz. Algebraische Theorie der Körpern, J. de Crelle, v. CXXXVII (1910), pp.~167-309 (new ed.~H. Hasse und R. Baer, Berlin-Leipzig (de Gruyter), 1930).

294b E. Steinitz. Rechteckige Systeme und Moduln in algebraischen Zahlkörpern, Math. Ann., v. LXXI (1912), pp.~328-354 与 LXXII (1912), pp.~.297-345.

295 S. Stevin. Les Oeuvres mathématiques . . . , ed.~A. Girard, Leyde (Elzevir), 1634.

296 E. Stiefel. Ueber eine Beziehung zwischen geschlossenen Lie'sche Gruppen und diskontinuierlichen Bewegungsgruppen euklidischer Räume und ihre Anwendung auf die Aufzählung der einfachen Lie'sche Gruppen, Comm. Math. Helv., v. XIV (1941-42), pp.~350-380.

297 T. Stieltjes. Recherches sur les fractions continues, Ann. Fac. Sci. de Toulouse, v. VIII (1894), pp.~J.1-J.122.

298 M. Stifel. Arithmetica integra, Nüremberg, 1544.

299 J. Stirling. Lineae tertii ordinis Newtonianae \ldots, Londini, 1717 (new ed., Paris (Duprat), 1797).

300 A. H. Stone. 仿紧性与乘积空间, Bull. Amer. Math. Soc., v. LIV (1948), pp.~977-982.

301a M. H. Stone. 布尔代数的表示理论, Trans. Amer. Math. Soc., v. XL (1936), pp.~37-111.

301b M. H. Stone. 布尔环理论对一般拓扑的应用, Trans. Amer. Math. Soc., v. XLI (1937), pp.~375-481.

301c M. H. Stone. 广义魏尔斯特拉斯逼近定理, Math. Magazine, v. XXI (1948), pp.~167-183 与 237-254.

302a C. Sturm. Sur les équations différentielles linéaires du second ordre, Journ. de Math., (1), v. I (1836), pp.~106-186.

302b C. Sturm. Sur une classe d'équations à différences partielles, Journ. de Math., (1), v. I (1836), pp.~373-444.

303 L. Sylow. Théorèmes sur les groupes de substitutions, Math. Ann., v. V (1872), pp.~584-594.

304 J. J. Sylvester. 数学论文集, 4 vol., Cambridge, 1904-1911.

305 B. Taylor. Methodus Incrementorum directa et inversa, Londini, 1715.

306 P. Tchebychef. Sur deux théorèmes relatifs aux probabilités, Acta Math., v. XIV (1890), pp.~305-315 (= Oeuvres, v. II, pp.~481-491).

307 H. Tietze. Über Funktionen die auf einer abgeschlossenen Menge stetig sind, J. de Crelle, v. CXLV (1915), pp.~9-14.

308a J. Tits. Groupes simple et géométries associées, Proc. Intern. Congress Math., Stockholm, 1962, pp.~197-221.

308b J. Tits. Théorème de Bruhat et sous-groupes paraboliques, C. R. Acad. Sci., v. CCLIV (1962), pp.~2910-2912.

308c J. Tits. 代数单群与抽象单群, 《数学年鉴》, v. LXXX (1964), pp.~313-329.

309a O. Toeplitz. Ueber die Auflösung unendlichvieler linearer Gleichungen mit unendlichvielen Unbekannten, Rend. Circ. Mat. Palermo, v. XXVIII (1909), pp.~88-96.

309b O. Toeplitz. Das Verhältnis von Mathematik und Ideenlehre bei Plato, Quellen und Studien zur Geschichte der Math., Abt. B: Studien, v. I (1929), pp.~3-33.

309c O. Toeplitz. Die mathematische Epinomisstelle, Quellen und Studien zur Geschichte der Math., Abt. B: Studien, v. II (1933), pp.~334-346.

310 E. Torricelli. Opere, 4 vol., ed G. Loria 与 G. Vassura, Faenza (Montanari), 1919.

311 J. Tropfke. Geschichte der Elementar-Mathematik, 7 vol., 2nd ed., Berlin-Leipzig (de Gruyter), 1921-24.

312 N. Tschebotarów. Grundzüge der Galois'schen Theorie (transl. Schwerdtfeger), Groningen (Noordhoff), 1950.

313 A. Tychonoff. Über die topologische Erweiterung von Räumen, Math. Ann., v. CII (1930), pp.~544-561.

314 P. Urysohn. Ueber die Mächtigkeit der zusammenhängenden Mengen, Math. Ann., v. XCIV (1925), pp.~262-295.

315 A. Vandermonde. Mémoire sur la résolution des équations, Hist. de l'Acad. royale des sciences, 1771 年, Paris (1774), pp.~365-416.

316 V. S. Varadarajan. 拓扑空间上的测度（俄文）, Mat. Sbornik, v. LV (1961), pp.~35-100（英文译本：Amer. Math. Soc. Transl. (2), vol.~48, pp.~161-228）.

317a B. L. van der Waerden. Moderne Algebra, 1st ed., 2 vol., Berlin (Springer), 1930-31.

317b B. L. van der Waerden. Die Klassification der einfachen Lieschen Gruppen, Math. Zeitschr., v. XXXVII (1933), pp.~446-462.

317c B. L. van der Waerden. Zenon und die Grundlagenkrise \ldots, Math. Ann., v. CXVII (1940), pp.~141-161.

317d B. L. van der Waerden. Die Arithmetik der Pythagoreer, I, Math. Ann., v. CXX (1947), pp.~127-153.

318 G. Veronese. Fondamenti di geometria, Padova, 1891.

319 Francisci Vietae. Opera mathematica \ldots, Lugduni Batavorum (Elzevir), 1646.

320a G. Vitali. Una proprieta delle funzioni misurabili, R. Ist. Lombardo Sci. Lett. Rend., (2), v. XXXVIII (1905), pp.~599-603.

320b G. Vitali. Sui gruppi di punti e sulla funzioni di variabili reali, Rend. Acc. Sci. di Torino, v. XLIII (1908), pp.~229-236.

321 H. Vogt. Die Entdeckungsgeschichte des Irrationalen nach Plato und andere Quellen des 4. Jahrhunderts, Bibl. Math., (3), v. X (1909), pp.~97-155.

322a V. Volterra. Leçons sur les fonctions de lignes, Paris (Gauthier-Villars), 1913.

322b V. Volterra. 泛函理论, London-Glasgow (Blackie 与 sons), 1930.

323 K. von Fritz. 梅塔蓬图姆的希帕索斯对不可公度性的发现, 《数学年鉴》, (2), v. XLVI (1945), pp.~242-264.

324a J. von Neumann. Eine Axiomatisierung der Mengenlehre, J. de Crelle, v. CLIV (1925), pp.~219-240.

324b J. von Neumann. Zur Theorie der Darstellung kontinuierlicher Gruppen, Berliner Sitzungsber., 1927, pp.~76-90.

324c J. von Neumann. Die Axiomatisierung der Mengenlehre, Math. Zeitschr., v. XXVII (1928), pp.~669-752.

324d J. von Neumann. Zur Hilbertschen Beweistheorie, Math. Zeitschr., v. XXVI (1927), pp.~1-46.

324e J. von Neumann. Zum Haarschen Mass in topologischen Gruppen, Compos. Math., v. I (1934), pp.~106-114.

324f J. von Neumann. 哈尔测度的唯一性, Mat. Sbornik., v. I (= XLIII) (1936), pp.~721-734.

324g J. von Neumann. 论算子环 III, 《数学年鉴》 (2), v. XLI (1940), pp.~94-161.

325 K. G. V. von Staudt. Beiträge zur Geometrie der Lage, Nürnberg, 1856.

326a J. Wallis. Opera Mathematica, 3 vol., Oxoniae, 1693-95.

326b J. Wallis. Logarithmotechnia Nicolai Mercatoris \ldots, Phil. Trans., v. III (1668), pp.~753-759.

327a H. Weber. Untersuchungen über die allgemeine Grundlagen der Galois'schen Gleichungstheorie, Math. Ann., v. XLIII (1893), pp.~521-544.

327b H. Weber. Leopold Kronecker, Math. Ann., v. XLIII (1893), pp.~1-25.

327c H. Weber. Beweis des Satzes, dass jede eigentlich primitive quadratische Form unendlich viele Primzahlen darzustellen fähig is, Math. Ann., v. XX (1882), pp.~301-329.

327d H. Weber. Theorie der Abels'chen Zahlkörper, IV, Acta Math., v. IX (1886), pp.~105-130.

328a J. Maclagan Wedderburn. 有限代数的一个定理, Trans. Amer. Math. Soc., v. VI (1905), pp.~349-352.

328b J. Maclagan Wedderburn. 论超复数, Proc. Lond. Math. Soc. (2), v. VI (1908), pp.~77-118.

328c J. Maclagan Wedderburn. 一类原始代数, Trans. Amer. Math. Soc., v. XV (1914), pp.~162-166.

328d J. Maclagan Wedderburn. 论可除代数, Trans. Amer. Math. Soc., v. XXII (1921), pp.~129-135.

329a K. Weierstrass. Mathematische Werke, 7 vol., Berlin (Mayer und Müller), 1894-1927.

329b K. Weierstrass. Briefe an P. du Bois-Reymond, Acta Math., v. XXXIX (1923), pp.~199-225.

330a A. Weil. Sur les fonctions elliptiques p-adiques, C. R. Acad. Sci., v. CCIII (1936), p.~22.

330b A. Weil. Sur les espaces uniformes et sur la topologie générale (Actual. Scient. et Ind., no 551), Paris (Hermann), 1937.

330c A. Weil. L'intégration dans les groupes topologiques et ses applications (Actual. Scient. et Ind., no 869), Paris (Hermann), 1940.

330d A. Weil. 代数几何基础 (Amer. Math. Soc. Coll. Publ., v. 29), new York, 1946.

330e A. Weil. 代数几何中的纤维空间（A. Wallace 记录的讲义）, Chicago Univ., 1952.

331a H. Weyl. Lettre à I. Schur, Berliner Sitzungsber., 1924, pp.~338-343.

331b H. Weyl. Theorie der Darstellung kontinuierlicher halbeinfacher Gruppen durch lineare Transformationen, Math. Zeitschr., v. XXIII (1925), pp.~271-309, v. XXIV (1926), pp.~328-395 与 789-791.

331c H. Weyl. 对称，Princeton (Princeton Univ. Press), 1952.

332 H. Weyl und F. Peter. Die Vollständigkeit der primitiven Darstellungen einer geschlossenen kontinuierlichen Gruppe, Math. Ann., v. XCVII (1927), pp.~737-755.

333 A. N. Whitehead. 论基数, 《美国数学杂志》, v. XXIV (1902), pp.~367-394.

334a J. H. C. Whitehead. 论无穷小群的分解, Proc. Camb. Phil. Soc., v. XXXII (1936), pp.~229-237.

334b J. H. C. Whitehead. 半单无穷小群的代数中的某些方程, 《数学季刊》, (2), v. VIII (1937), pp.~220-237.

335 H. Wieleitner. Der ``Tractatus de latitudinibus formarum'' des Oresme, Bibl. Math. (3), v. XIII (1912), pp.~115-145.

336 N. Wiener. 微分空间，J. Math. Phys. M. I. T., v. II (1923), pp.~131-174.

337a E. Witt. Theorie der quadratischen Formen in beliebigen Körpern, J. de Crelle, v. CLXXVI (1937), pp.~31-44.

337b E. Witt. Treue Darstellung Lieschen Ringe, J. de Crelle, v. CLXXVII (1937), pp.~152-160.

337c E. Witt. Spiegelungsgruppen und Aufzählung halbeinfacher Liescher Ringe, Abh. math. Sem. Univ. Hamburg, v. XIV (1941), pp.~289-322.

338 E. Wright. 纬度表, 1599.

339a W. H. Young. 论上积分与下积分, Proc. London Math. Soc. (2), v. II (1905), pp.~52-66.

339b W. H. Young. 积分理论中的一个新方法, Proc. London Math. Soc. (2), v. IX (1911), pp.~15-50.

340a O. Zariski. 抽象代数函数域的黎曼流形的紧性, Bull. Amer. Math. Soc., v. L (1944), pp.~683-691.

340b O. Zariski. 广义局部环, Summa Bras. Math., v. I (1946), pp.~169-195.

341 H. Zassenhaus. Lehrbuch der Gruppentheorie, Bd. I, Leipzig-Berlin (Teubner), 1937.

342a E. Zermelo. Beweis dass jede Menge wohlgeordnet werden kann, Math. Ann., v. LIX (1904), pp.~514-516.

342b E. Zermelo. Neuer Beweis für die Möglichkeit einer Wohlordnung, Math. Ann., v. LXV (1908), pp.~107-128.

342c E. Zermelo. Untersuchung über die Grundlagen der Mengenlehre, Math. Ann., v. LXV (1908), pp.~261-281.

343 G. Zolotareff. Sur la théorie des nombres complexes, Journ. de Math. (3), v. VI (1880), pp.~51-84 与 129-166.

344 M. Zorn. 关于超限代数中方法的一点注记, Bull. Amer. Math. Soc., v. XLI (1935), pp.~667-670.

345 A. Zygmund. 三角级数, 2 vol., Cambridge (Univ. Press), 1959.
